% Generated draft: requires mathematical and source review.
\documentclass[11pt]{article}
\usepackage[margin=25mm]{geometry}
\usepackage{fontspec}
\setmainfont{Latin Modern Roman}
\usepackage{amsmath,amssymb,mathtools,unicode-math}
\setmathfont{Latin Modern Math}
\usepackage{enumitem,needspace,xurl,hyperref}
\setlength{\parindent}{0pt}
\setlength{\parskip}{5pt}
\setlength{\emergencystretch}{3em}
\newcommand{\N}{\mathbb{N}}
\newcommand{\Z}{\mathbb{Z}}
\newcommand{\Q}{\mathbb{Q}}
\newcommand{\R}{\mathbb{R}}
\newcommand{\C}{\mathbb{C}}
\newcommand{\F}{\mathbb{F}}
\newcommand{\A}{\mathbb{A}}
\newcommand{\PP}{\mathbb{P}}
\newcommand{\E}{\mathbb{E}}
\newcommand{\eps}{\varepsilon}
\newcommand{\dd}{\,\mathrm{d}}
\begin{document}
\section*{UnsolvedMath ID: 30003577}
\textbf{Source problem code:} catalogue entry
\subsection{Definitions and mathematical statement}
Under the conventional interpretation of the modular surface, let
\[
X=\operatorname{PSL}_2(\Z)\backslash\mathbb H,
\qquad \mathbb H=\{x+iy:y>0\},
\qquad d\mu=\frac{3}{\pi}\frac{dx\,dy}{y^2},
\]
so that \(\mu(X)=1\). Use the nonnegative Laplacian
\[
\Delta=-y^2(\partial_x^2+\partial_y^2).
\]
A weight-zero Hecke--Maass cusp form \(\phi\) is a smooth automorphic function on \(\mathbb H\), square integrable on \(X\), that is a simultaneous eigenfunction of \(\Delta\) and the usual Hecke operators and has vanishing constant Fourier coefficient at the cusp. Thus
\[
\Delta\phi=\lambda_\phi\phi,
\qquad
\int_0^1\phi(x+iy)\,dx=0\quad(y>0).
\]
Normalize \(\int_X|\phi|^2\,d\mu=1\), and define the probability measure
\[
d\mu_\phi=|\phi|^2\,d\mu.
\]
These specify standard conventions for the surface, Laplacian, and normalization; the supplied record does not explicitly give them.

\textbf{Requested assertion (strong arithmetic quantum unique ergodicity).} Establish equidistribution of \(\mu_\phi\) toward \(\mu\) with a power saving as \(\lambda_\phi\to\infty\) through Hecke--Maass cusp forms. A precise formulation using fixed smooth compactly supported test functions would be
\[
\exists\delta>0\ \ \forall a\in C_c^\infty(X)\ \ \exists C_a>0\ \ \exists L_a>0\ \ \forall\phi:
\quad
\lambda_\phi\ge L_a
\ \Longrightarrow\ 
\left|\int_X a\,d\mu_\phi-\int_X a\,d\mu\right|
\le C_a\lambda_\phi^{-\delta},
\]
where \(\phi\) ranges over all forms satisfying the conditions above, and \(\delta\) is independent of both \(a\) and \(\phi\).

\textbf{Unresolved scope.} The source specifies a rate \(O(\lambda^{-\delta})\) but does not specify a topology or testing class for comparing measures, the permitted dependence of the implied constant, or whether a common exponent is required for all test functions. Consequently, the displayed quantified inequality is an explicitly proposed interpretation of the requested rate, not a uniquely determined transcription. In particular, it supplies no uniform estimate for varying test functions or sets and no separate quantitative estimate in the cusp. The intended scope requires human mathematical review.
\subsection{Short English statement}
For normalized Hecke--Maass cusp forms on the modular surface, the problem asks for their probability densities to approach normalized hyperbolic area at a power rate in the Laplace eigenvalue. A proposed precise interpretation tests against each fixed smooth compactly supported function, with one positive exponent valid for all forms and tests. The supplied record leaves the testing class and uniformity requirements unspecified.
\subsection{Sources}
\begin{itemize}
\item Local UnsolvedMath registry: supplied complete record with permanent ID \(30003577\), title \emph{Power-Rate Arithmetic Quantum Unique Ergodicity}, source problem code \texttt{catalogue entry}, in the Oberwolfach Reports Open Problems collection. No individual proposer is supplied; the proposed year is \(2017\).
\item Supplied report citation: \emph{Automorphic Forms and Arithmetic (2018)}, Oberwolfach Reports, DOI \texttt{10.4171/owr/2017/40}, pp.\ 2475--2537. Supplied source URL and literature-source link labelled ``OWR report'': \url{https://doi.org/10.4171/owr/2017/40}. The record describes Oberwolfach Reports as documenting mathematical workshops and their open-problem discussions.
\item Supplied literature-source link labelled ``Nelson'': \url{https://arxiv.org/abs/1702.02908}. No further bibliographic details are supplied.
\item Literature assessment dated exactly \texttt{2026-08-21}: the supplied curation labels the problem open and states that arithmetic quantum unique ergodicity is known, including level-aspect results, while the power-saving rate asserted by strong AQUE remains open. This is supplied dated curation derived from the OWR material, not an independently verified literature assessment. The record expressly requests independent verification of time-sensitive status and cited claims.
\item The supplied statement-status label is \texttt{corrected\_verified}; its accompanying explanation concerns preservation of mathematical content while repairing notation, typography, and sentence structure for publication. This editorial label does not establish mathematical correctness or resolve the formulation issues above. The report citation's \(2018\) designation and its DOI containing \(2017\), together with the supplied proposed year \(2017\), are retained without resolving their bibliographic relationship.
\end{itemize}
\end{document}
